\section*{第 \ref{ch:b2:angiosperm-reproduction} 章——被子植物的有性生殖}

\begin{solution}{exo:b2:angiosperm-reproduction:1}
萼片（花萼）、花瓣（花冠）、\omterm{def:b2:angiosperm-reproduction:flower}{雄蕊}（\omterm{def:b2:angiosperm-reproduction:flower}{雄蕊}群）、\omterm{def:b2:angiosperm-reproduction:flower}{心皮}
（雌蕊群）。\omterm{def:b2:angiosperm-reproduction:flower}{雄蕊}：花丝和花药。\omterm{def:b2:angiosperm-reproduction:flower}{心皮}：\omterm{def:b2:angiosperm-reproduction:flower}{柱头}、花柱、
\omterm{def:b2:angiosperm-reproduction:flower}{子房}（含\omterm{def:b2:plant-life-cycles:heterospory}{胚珠}）。以上均为\omterm{def:b2:meiosis-heredity:meiosis}{二倍体}\omterm{def:b2:plant-life-cycles:cycles}{孢子体}，只有\omterm{prop:b2:angiosperm-reproduction:pollen}{花粉粒}
（位于花药内）和\omterm{prop:b2:angiosperm-reproduction:embryosac}{胚囊}（位于\omterm{def:b2:plant-life-cycles:heterospory}{胚珠}内）是
\omterm{def:b2:meiosis-heredity:meiosis}{单倍体}\omterm{def:b2:plant-life-cycles:cycles}{配子体}。
\end{solution}

\begin{solution}{exo:b2:angiosperm-reproduction:2}
\omterm{def:b2:plant-life-cycles:heterospory}{小孢子} $n$；营养细胞 $n$；精子 $n$；卵 $n$；\omterm{prop:b2:angiosperm-reproduction:embryosac}{极
核} $n$；受精卵 $2n$；\omterm{prop:b2:angiosperm-reproduction:double}{胚乳} $3n$；\omterm{def:b2:angiosperm-reproduction:seed}{种皮} $2n$（母体珠被）；
果皮 $2n$（母体\omterm{def:b2:angiosperm-reproduction:flower}{子房}）。
\end{solution}

\begin{solution}{exo:b2:angiosperm-reproduction:3}
风媒植物（禾草、栎、桦、榛、车前）：花小而绿，无
花瓣、香味或\omterm{def:b2:angiosperm-reproduction:pollination}{花蜜}，花药悬垂于长花丝上，
\omterm{def:b2:angiosperm-reproduction:flower}{柱头}呈羽毛状，\omterm{def:b2:plant-life-cycles:heterospory}{花粉}量大、光滑干燥，在叶片展开
之前开花。蜂媒植物（三叶草、鼠尾草、毛地黄、薰衣草、苹果）：花瓣有色，
常具紫外线标记、着陆平台及香味，
分泌\omterm{def:b2:angiosperm-reproduction:pollination}{花蜜}，\omterm{def:b2:plant-life-cycles:heterospory}{花粉}量中等、带黏性且表面有纹饰，在
白天开花。
\end{solution}

\begin{solution}{exo:b2:angiosperm-reproduction:4}
同一\omterm{thm:b2:angiosperm-reproduction:tube}{花粉管}的两个精子融合，一个与卵
融合形成受精卵（$2n$，即胚），另一个与中央细胞
融合形成\omterm{prop:b2:angiosperm-reproduction:double}{胚乳}核（$3n$，即养料贮存结构）。\omterm{prop:b2:angiosperm-reproduction:double}{胚乳}只在受精后才开始发育，
因此植物只为含胚的\omterm{def:b2:plant-life-cycles:heterospory}{胚珠}供给养料，
这与裸子植物不同，后者在受精前即已贮备养料。
\end{solution}

\begin{solution}{exo:b2:angiosperm-reproduction:5}
$25/1.2 = \qty{21}{h}$。第 5、6、7 天伸出的花丝可接受\omterm{def:b2:plant-life-cycles:heterospory}{花粉}；
第 8 至 12 天伸出的则不能，因此 $3/8$ 的花丝获得授粉，
果穗有八分之五不结实。
\end{solution}

\begin{solution}{exo:b2:angiosperm-reproduction:6}
$S_1S_2\times S_1S_2$：0。$S_1S_2\times S_1S_3$：$\tfrac12$（仅 $S_3$
\omterm{def:b2:plant-life-cycles:heterospory}{花粉}可萌发）。$S_1S_2\times S_3S_4$：全部亲和。第二个杂交的后代：
卵为 $S_1$ 或 $S_2$，精子为 $S_3$：$S_1S_3$ 和 $S_2S_3$，各
占一半。
\end{solution}

\begin{solution}{exo:b2:angiosperm-reproduction:7}
一株植物携带 $n$ 个\omterm{def:b2:meiosis-heredity:vocabulary}{等位基因}中的 2 个，并拒绝携带
其中任一者的\omterm{def:b2:plant-life-cycles:heterospory}{花粉}。亲和比例为 $(n - 2)/n$。新\omterm{def:b2:meiosis-heredity:vocabulary}{等位基因}
不会被任何雌蕊拒绝，因此其\omterm{def:b2:plant-life-cycles:heterospory}{花粉}的结实数多于平均水平，该\omterm{def:b2:meiosis-heredity:vocabulary}{等位基因}
会扩散，直至与其他\omterm{def:b2:meiosis-heredity:vocabulary}{等位基因}同样常见；同样的优势保护每一个
稀有\omterm{def:b2:meiosis-heredity:vocabulary}{等位基因}不致丢失。因此自然选择会使\omterm{def:b2:meiosis-heredity:vocabulary}{等位基因}不断积累，
\omterm{prop:b2:angiosperm-reproduction:selfing}{自交不亲和}物种的自然种群往往携带数十个。
\end{solution}

\begin{solution}{exo:b2:angiosperm-reproduction:8}
两个\omterm{prop:b2:angiosperm-reproduction:embryosac}{极核}来自同一\omterm{def:b2:plant-life-cycles:heterospory}{大孢子}，因此中央细胞
为 $YY$ 或 $yy$（各占一半）；精子为 $Y$ 或 $y$（各
占一半）：\omterm{prop:b2:angiosperm-reproduction:double}{胚乳}为 $YYY$、$YYy$、$Yyy$、$yyy$，各占 $\tfrac14$，故
四分之三为黄色。由 $yy$ 授粉时：$YYy$ 和 $yyy$，各占一半——
一半籽粒为黄色。
\end{solution}

\begin{solution}{exo:b2:angiosperm-reproduction:9}
将大麦籽粒纵切为两半；含胚的一半消化
其淀粉，无胚的一半则不消化，除非
加入赤霉素，此时它们才会消化。(a) 胚是信号的来源，而不是
酶的来源。(b) 赤霉素即是信号，本身即足以
起作用。(c) 淀粉酶由\omterm{prop:b2:angiosperm-reproduction:double}{胚乳}的糊粉层合成，
糊粉层对该激素作出响应。无胚的一半将信号与
响应分开：若没有这一半，便无法判断胚
是否自身消化淀粉。
\end{solution}

\begin{solution}{exo:b2:angiosperm-reproduction:10}
风媒：$10^{6}\times\qty{1}{\micro g} = \qty{1}{g}$ 每个\omterm{def:b2:plant-life-cycles:heterospory}{胚珠}。动物媒：
\qty{1}{mg} \omterm{def:b2:plant-life-cycles:heterospory}{花粉}与 \qty{0.5}{mg} 糖类合计 $= \qty{1.5}{mg}$
每个\omterm{def:b2:plant-life-cycles:heterospory}{胚珠}——动物媒约便宜七百倍。风媒
无需伙伴：它可在早春昆虫飞行之前进行，在
寒冷、多风和空旷的地方，夜间与雨中，在单一物种
的密集群落中进行，并且不会被\omterm{def:b2:angiosperm-reproduction:pollination}{花蜜}盗贼欺骗，也不会
因\omterm{def:b2:angiosperm-reproduction:pollination}{传粉}者灭绝而被遗弃。
\end{solution}

\begin{solution}{exo:b2:angiosperm-reproduction:11}
落在母株基部的\omterm{def:b2:angiosperm-reproduction:seed}{种子}落于其荫蔽之下，位于其
根之间，也是其专一的\omterm{def:b2:angiosperm-reproduction:seed}{种子}捕食者和病原体
聚集之处；几乎全部死亡。鸟类可将\omterm{def:b2:angiosperm-reproduction:seed}{种子}带出数公里
（\qty{12}{km}，仅需二十分钟），连同肥料一起丢在
树篱或空地中；即便一百粒中只有一粒落在适宜地点，
其数量也多于树下的全部\omterm{def:b2:angiosperm-reproduction:seed}{种子}，并且种群扩散的
速度取决于鸟类，而非落果。果肉是
这张车票的代价。
\end{solution}

\begin{solution}{exo:b2:angiosperm-reproduction:12}
包被使\omterm{def:b2:plant-life-cycles:heterospory}{花粉}在母体组织中生长，从而使雌蕊能够
检验\omterm{def:b2:plant-life-cycles:heterospory}{花粉}（\omterm{prop:b2:angiosperm-reproduction:selfing}{自交不亲和}、排斥异种），使许多\omterm{thm:b2:angiosperm-reproduction:tube}{花粉管}
竞争，最快的一个使\omterm{def:b2:plant-life-cycles:heterospory}{胚珠}受精；同时保护\omterm{def:b2:plant-life-cycles:heterospory}{胚珠}免受干燥
和食草动物的侵害，并使\omterm{def:b2:angiosperm-reproduction:flower}{子房}
转变为吸引动物传播\omterm{def:b2:angiosperm-reproduction:seed}{种子}的\omterm{def:b2:angiosperm-reproduction:seed}{果实}；\omterm{prop:b2:angiosperm-reproduction:double}{双受精}
则只为已受精的\omterm{def:b2:plant-life-cycles:heterospory}{胚珠}提供养料。代价是：\omterm{prop:b2:angiosperm-reproduction:pollen}{花粉粒}
必须到达\omterm{def:b2:angiosperm-reproduction:flower}{柱头}而非\omterm{def:b2:plant-life-cycles:heterospory}{胚珠}本身，因此需要
花柱、\omterm{thm:b2:angiosperm-reproduction:tube}{花粉管}及其导向机制，且\omterm{def:b2:angiosperm-reproduction:seed}{果实}
是一笔沉重的投资。这一权衡如此有利于被子植物，
它们在一亿
年间由几乎为零增至占植物物种的九成。
\end{solution}

\begin{solution}{pb:b2:angiosperm-reproduction:1}
\textbf{1.} 对 A（$S_1S_2$）：\omterm{def:b2:plant-life-cycles:heterospory}{花粉} A 的 $f = 0$；B（$S_1$、$S_3$）$f =
\tfrac12$；C（$S_3$、$S_4$）$f = 1$。
\textbf{2.} 对 B（$S_1S_3$）：A 为 $\tfrac12$，B 为 0，C 为 $\tfrac12$。对 C
（$S_3S_4$）：A 为 1，B 为 $\tfrac12$，C 为 0。
\textbf{3.} 50 粒可亲和\omterm{def:b2:plant-life-cycles:heterospory}{花粉}；$0.1\times 50 = 5$ 个\omterm{def:b2:plant-life-cycles:heterospory}{胚珠}。
\textbf{4.} 来自 C 的有 100 粒可亲和\omterm{def:b2:plant-life-cycles:heterospory}{花粉}，$0.1\times 100 = 10$，即全部
10 个\omterm{def:b2:plant-life-cycles:heterospory}{胚珠}；来自 A 的则无。
\textbf{5.} 来自 C 访花所结的\omterm{def:b2:angiosperm-reproduction:seed}{果实}（10 粒\omterm{def:b2:angiosperm-reproduction:seed}{种子}）保留；来自 B 的（5 粒\omterm{def:b2:angiosperm-reproduction:seed}{种子}）
同样保留；来自 A 的则被淘汰。仅由 A 组成的果园不结果。
\textbf{6.} 一半的花结实（即来自 C 的访花）：每棵树 2500 个
\omterm{def:b2:angiosperm-reproduction:seed}{果实}，是所需 250 个的十倍——足额丰收，而
树会脱落多余的\omterm{def:b2:angiosperm-reproduction:seed}{果实}。
\textbf{7.} 它在各品种的整个花期内都散出可亲和的\omterm{def:b2:plant-life-cycles:heterospory}{花粉}；
但若它与某品种共享两个 $S$ \omterm{def:b2:meiosis-heredity:vocabulary}{等位基因}，
则在该品种上无法授粉。
\textbf{8.} $8\times 10^{7}/(7\times 86400) = 132$ 粒每平方
米每秒。
\textbf{9.} $132/0.25 = 530$ 粒每立方米。
\textbf{10.} $132\times 10^{-4} = 0.013$ 粒每秒：每
\qty{76}{s} 一粒。
\textbf{11.} $0.013\times 86400 \approx 1100$ 粒每天；最先
到达\omterm{def:b2:plant-life-cycles:heterospory}{胚珠}的\omterm{thm:b2:angiosperm-reproduction:tube}{花粉管}使其受精，此后\omterm{def:b2:plant-life-cycles:heterospory}{胚珠}
即对其余\omterm{thm:b2:angiosperm-reproduction:tube}{花粉管}关闭。
\textbf{12.} 授粉后 \qty{20}{h}。
\textbf{13.} \omterm{def:b2:plant-life-cycles:heterospory}{花粉}在第 0--7 天散出，花丝自第 5 天开始伸出：只有第 5、6 和
7 天重叠——占花丝的 $3/8$，而在第 7 天之后伸出的花丝
得不到\omterm{def:b2:plant-life-cycles:heterospory}{花粉}；果穗大多空秕。吐丝期干旱是
玉米歉收的典型原因。
\textbf{14.} 胚为 $Yy$；\omterm{prop:b2:angiosperm-reproduction:double}{胚乳}为 $YYy$。
\textbf{15.} 黄色，因为 $Y$ 为\omterm{def:b2:meiosis-heredity:vocabulary}{显性}：父本的性状
在母本植株所结的\omterm{def:b2:angiosperm-reproduction:seed}{种子}中显现——即\omterm{def:b2:plant-life-cycles:heterospory}{花粉}直感现象。
\textbf{16.} 中央细胞为 $YY$ 或 $yy$：\omterm{prop:b2:angiosperm-reproduction:double}{胚乳} $YYy$（占一半）和
$yyy$（占一半）。
\textbf{17.} $YYy$：20 个单位；$Yyy$：10 个单位；
一个 $yyY$（其中的 $Y$ 来自父本）与 $Yyy$ 为同一\omterm{def:b2:meiosis-heredity:vocabulary}{基因型}，同为 10 个单位。
不可能存在两份父本剂量：一个精子使中央细胞受精，而
\omterm{prop:b2:angiosperm-reproduction:double}{胚乳}中母本与父本基因组之比始终为二比一。
\textbf{18.} 它是淀粉贮藏库；$0.3/0.33 = \qty{91}{\%}$ 的
籽粒。
\textbf{19.} 贮藏库仅在受精后才形成，因此不会为
不含胚的\omterm{def:b2:plant-life-cycles:heterospory}{胚珠}投入资源；裸子植物的\omterm{def:b2:plant-life-cycles:cycles}{配子体}
则预先形成，若授粉失败便被浪费。
\textbf{20.} $0.9\times 600 = 540$ 粒每穗、每株；
$4320$ 粒每平方米；$4320\times 0.3 = \qty{1.3}{kg}$ 的\omterm{prop:b2:angiosperm-reproduction:double}{胚乳}
每平方米。
\textbf{21.} \qty{13}{t/ha}。
\textbf{22.} 籽粒碳 $0.5\times 1.2 = \qty{0.6}{kg/m^2}$；
\omterm{prop:b2:angiosperm-reproduction:double}{胚乳}碳 $0.45\times 1.3 = \qty{0.58}{kg/m^2}$：相符
（余下的少量为胚和\omterm{def:b2:angiosperm-reproduction:seed}{种皮}）。
\textbf{23.} $8\times 10^{7}/4320 \approx \num{18500}$ 粒\omterm{def:b2:plant-life-cycles:heterospory}{花粉}
每籽粒。
\textbf{24.} 其自身的\omterm{def:b2:plant-life-cycles:heterospory}{花粉}云稀薄且被风吹散，其
雄穗大多在自身花丝伸出之前已散粉，因此许多花丝
未能授粉，果穗便出现缺粒。
\textbf{25.} A 田块：一半的花结实，每棵树 2500 个，足额
丰收。田间：每平方米 4320 粒籽粒，\qty{13}{t/ha}。
\end{solution}
